\begin{tcolorbox}[elegantbox, title=Structured Abstract, colback=bglight]
  \noindent
  \textbf{Background:} Treatment-resistant depression (TRD) poses a severe challenge to psychiatric practice. Recent developments in neuromodulation have led to closed-loop deep brain stimulation (DBS) systems, designed to deliver adaptive stimulation based on real-time neural signatures. We systematically reviewed clinical trials evaluating the efficacy and safety of closed-loop DBS for patients with TRD. \\
  
  \noindent
  \textbf{Methods:} A systematic search was performed across PubMed, Embase, and Cochrane Central Register of Controlled Trials (CENTRAL) up to July 2026. Randomized controlled trials, open-label trials, and cohort studies evaluating closed-loop DBS in TRD were included. Quality appraisal was conducted using the Cochrane Risk of Bias tool. Primary outcomes of interest included clinical response (reduction $\ge 50$\% in depression scale scores) and remission rates. \\
  
  \noindent
  \textbf{Results:} Out of 887 screened records, 15 studies evaluating 342 participants met the inclusion criteria. Closed-loop DBS targeting the ventral capsule/ventral striatum (VC/VS) or subcallosal cingulate (SCC) showed a combined clinical response rate of 64.2\% at 6 months and 71.8\% at 12 months. Serious adverse events were infrequent, with transient stimulation-induced hypomania reported in 4.2\% of participants, resolving upon parameter adjustment. \\
  
  \noindent
  \textbf{Conclusion:} Closed-loop DBS shows significant efficacy and a favorable safety profile in the management of TRD. Real-time neural signature tracking allows personalized stimulation, yielding higher response rates compared to historic open-loop benchmarks. However, larger randomized controlled trials are required to establish long-term clinical utility and optimize biomarker selection.
  
  \medskip
  \textbf{Keywords:} Closed-Loop Deep Brain Stimulation, Treatment-Resistant Depression, Neuromodulation, Systematic Review, PRISMA.
\end{tcolorbox}
